\documentclass{article}
\usepackage{bm}
\usepackage{braket}
\usepackage{lipsum}
\usepackage{mathtools}
\usepackage{amsthm}
\usepackage{fancyhdr}
\usepackage{float}
\usepackage{listings}
\lstset{breaklines=true}
\usepackage{graphicx}
\usepackage{xcolor}
\usepackage{titlesec}
\usepackage{titling}
\usepackage{tabularx}
\usepackage{enumitem}
\usepackage{amsmath}
\usepackage{amsfonts}
\usepackage{hyperref}
\usepackage[margin=1in]{geometry}

\hypersetup{
	colorlinks=true,
	linkcolor=blue,
	filecolor=magenta,
	urlcolor=blue,
	citecolor=blue,
	anchorcolor=blue
}

\renewcommand{\thesubsubsection}{\Roman{subsubsection}}
\title{Math 3IA3 - Midterm}
\author{Matthew Farah, \href{mailto:farahm11@mcmaster.ca}{$\langle$farahm11@mcmaster.ca$\rangle$}, 400468251}
\date{\today}

\makeatletter

\pagestyle{fancy}
\fancyhead{}
\fancyhead[L]{Matthew Farah, $\langle$\href{mailto:farahm11@mcmaster.ca}{farahm11@mcmaster.ca}$\rangle$, 400468251}
\fancyhead[R]{Math 3IA3 - Midterm} 

\newcolumntype{Y}{>{\centering\arraybackslash}X}
\setlength{\parindent}{0cm}

\newcounter{problem}
\renewcommand{\theproblem}{\arabic{problem}}

\newenvironment{problem}{
	\newpage
	\refstepcounter{problem}
	\addcontentsline{toc}{section}{\protect{\large{Problem \theproblem}}}
	\par
	\large\textbf{Problem \theproblem}
	\vspace{.2cm}
	\par
	\setcounter{step}{0}
}{\vspace{.5em}}

\newenvironment{solution}{
	\vspace{.5em}\par\large\textbf{{Solution:}}\par\vspace{.5em}
}{
	\vspace{.5em}
}

\newcounter{step}
\renewcommand{\thestep}{\Roman{step}}

\newenvironment{step}{
	\refstepcounter{step}
	\vspace{1.5em}\par\noindent\large\textbf{(\thestep)}
	\setcounter{substep}{0}
}{
	\par
	\vspace{1.5em}
}

\makeatother

\newcounter{substep}
\renewcommand{\thesubstep}{\alph{substep}}

\newenvironment{substep}{
	\refstepcounter{substep}
	\vspace{1.5em}\par\noindent\large\textbf{(\thesubstep)}
}{
	\par
	\vspace{1.5em}
}

\makeatother

\begin{document}

\maketitle
\pagestyle{fancy}

\tableofcontents
\newpage

\begin{problem}
	Let $S$ be a subset of $\mathbb{R}$ which is bounded below. Prove that a lower bound $w$ of $S$ is the infimum of $S$ if and only if for every $\epsilon > 0$, there exists some $t \in S$ such that $t < w + \epsilon$.
\end{problem}

\begin{solution}
	\begin{proof}
		$ $
		\begin{step}
			Prove the forward direction.
		\end{step}
	
		If $w = \inf{S}$, then $w \le s, \forall s \in S$ and for all $l \in \mathbb{R}$ such that $l$ lowerbounds $S$, $l \le w$ by definition of infimum. \\
	
		Now, let $\epsilon > 0$ be arbitrary, and assume, for the purposes of deriving a contradiction, that there exists no element $s_0 \in S$ such that $s_0 < w + \epsilon$. This therefore implies that $w + \epsilon$ must be a lower bound of $S$ since $\forall s \in S, w + \epsilon \le s$. However, clearly $w < w + \epsilon$, which implies that there exists a lower bound of $S$ which is larger than $w$. This contradicts our prior assumptions that $w$ must be the greatest lower bound of $S$, and thus it must be the case that for arbitrary $\epsilon > 0$, $\exists s \in S$ such that $s < w + \epsilon$.
	
		\begin{step}
			Prove the reverse direction.
		\end{step}
	
		If $w$ is a lower bound of $S$ with the property that $\forall \epsilon > 0$, $\exists s_0 \in S$ such that $s_0 < w + \epsilon$, this implies that $s_0 - \epsilon < w$. \\
	
		Let $l$ be an arbitrary lower bound of $S$. Then we know for $l$ to lower bound $S$, $l \le s_0$, however, this implies that $l < w + \epsilon$ from our previous deduction regarding the relationship between $s_0$ and $w$. This therefore implies that $l \le w$, which means that $w$, which is a lower bound of $S$ is larger than or equal to any other possible lower bound of $S$. Therefore, by definition of infimum, we know that $w = \inf{S}$.
	\end{proof}
\end{solution}

\begin{problem}
	If $0 < a < b$, determine $\lim_{n\to\infty}\frac{a^{n + 1} + b^{n + 1}}{a^n + b^n}$.
\end{problem}

\begin{solution}
	\begin{align*}
		\lim_{n\to\infty} \left( \frac{a^{n + 1} + b^{n + 1}}{a^n + b^n} \right) &= \lim_{n\to\infty} \left( \frac{b^{n} \left( \frac{a^{n + 1}}{b^n} + b \right)}{b^n \left( \frac{a^n}{b^n} + 1 \right)} \right) \\
		&= \lim_{n\to\infty} \left( \frac{\frac{a^{n + 1}}{b^n} + b}{\frac{a^n}{b^n} + 1} \right) \\
		&= \lim_{n\to\infty} \left( \frac{a \left( \frac{a}{b} \right)^n + b}{\left( \frac{a}{b} \right)^n + 1} \right) \\
		&= \frac{a \lim_{n\to\infty}\left( \frac{a}{b} \right)^n + b}{\lim_{n\to\infty} \left( \frac{a}{b} \right)^n + 1} \\
	\end{align*}

	Now, since $0 < a < b$, $\frac{a}{b} < 1$. Therefore, we know that $\lim_{n\to\infty}\left( \frac{a}{b} \right)^n = 0$. Thus, we can conclude that:

	\begin{align*}
		\frac{a \lim_{n\to\infty}\left( \frac{a}{b} \right)^n + b}{\lim_{n\to\infty} \left( \frac{a}{b} \right)^n + 1} &= \frac{a \cdot 0 + b}{0 + 1} \\
		&= \frac{0 + b}{1} \\
		&= \frac{b}{1} \\
		&= b \\
	\end{align*}

	Thus, $\lim_{n\to\infty} \left( \frac{a^{n + 1} + b^{n + 1}}{a^n + b^n} \right) = b$.
\end{solution}

\begin{problem}
	Let $x_1 = 1$ and let $x_{n + 1} = \sqrt{2 + x_n}, \forall n \in \mathbb{N}$ such that $n \ge 1$. Show that $\lim_{n\to\infty}(x_n)$ exists and find the limit.
\end{problem}

\begin{solution}
	\begin{step}
		Prove $(x_n)$ is upperbounded.
	\end{step}

	Selecting 10 as an arbitrary number I conjecture upperbounds the sequence $(x_n)$, I will attempt to prove that 10 does indeed upperbond $(x_n)$ by induction. Let $P(n)$ be the statement $P(n): x_n \le 10$.

	\begin{substep}
		Show the base case holds (Show $P(1)$ is true).
	\end{substep}

	\begin{align*}
		&\text{LHS:} & \text{RHS:} \\
		& = x_1 & = 10 \\
		& = 1 & = 10 \\
	\end{align*}

	Therefore, since the LHS $\le$ RHS, the base case holds and $P(1)$ is true.

	\begin{substep}
		Formulate the induction hypothesis (State $P(k)$).
	\end{substep}

	Let $k \in \mathbb{N}$ be an arbitrary natural number for which our statement holds true. We can therefore state our induction hypothesis to be:

	\begin{align*}
		P(k): x_k \le 10
	\end{align*}

	\begin{substep}
		Prove the induction step (Show $P(k) \implies P(k + 1)$).
	\end{substep}

	\begin{align*}
		x_{k + 1} &= \sqrt{2 + x_k} \qquad \langle \; \text{By definition of} \; (x_n) \; \rangle \\
		&\le \sqrt{2 + 10} \qquad \langle \; \text{From our induction hypothesis} \; \rangle \\
		&= \sqrt{12} \\
		&\le 10 \\
	\end{align*}

	Therefore, since $P(1)$ is true, and for arbitrary $k \in \mathbb{N}$ where $P(k)$ holds true, $P(k) \implies P(k + 1)$, then by the principle of mathematical induction, $P(n)$ is true $\forall n \in \mathbb{N}$ and thus 10 is an upper bound of $(x_n)$.

	\begin{step}
		Prove that $(x_n)$ is increasing.
	\end{step}

	By computing the first few values of $(x_n)$, I conjecture, and will attempt to prove by induction, that $(x_n)$ is an increasing sequence. For this purpose, let $P(n)$ be the statement $P(n): x_{n} \le x_{n + 1}$

	\begin{substep}
		Show the base case holds (Show $P(1)$ is true).
	\end{substep}

	\begin{align*}
		&\text{LHS:} & \text{RHS:} \\
		& = x_1 & = x_2 \\
		& = 1 & = \sqrt{2 + 1} \\
		& = 1 & = \sqrt{3} \\
	\end{align*}

	Therefore, since the LHS $\le$ RHS, the base case holds and $P(1)$ is true.

	\begin{substep}
		Formulate the induction hypothesis (State $P(k)$).
	\end{substep}

	Let $k \in \mathbb{N}$ be an arbitrary natural number for which our statement holds true. We can therefore state our induction hypothesis to be:

	\begin{align*}
		P(k): x_k \le x_{k + 1}
	\end{align*}

	\begin{substep}
		Prove the induction step (Show $P(k) \implies P(k + 1)$).
	\end{substep}

	\begin{align*}
		x_{k + 2} &= \sqrt{2 + x_{k + 1}} \qquad \langle \; \text{By definition of} \; (x_k) \; \rangle \\
		&\ge \sqrt{2 + x_k} \qquad \langle \; \text{By induction hypothesis} \; \rangle \\
		&= x_k
	\end{align*}

	Therefore, since $P(1)$ is true, and for arbitrary $k \in \mathbb{N}$ where $P(k)$ holds true, $P(k) \implies P(k + 1)$, then by the principle of mathematical induction, $P(n)$ is true $\forall n \in \mathbb{N}$ and thus $(x_n)$ is an increasing sequence. \\

	Thus, by $\textbf{(I)}$ and $\textbf{(II)}$, $(x_n)$ is an increasing sequence which is bounded above. Therefore, by the Monotone Convergence Theorem, $(x_n)$ converges and thus $\lim_{n\to\infty}(x_n)$ exists.

	\begin{step}
		Find $\lim_{n\to\infty}(x_n)$.
	\end{step}

	Let $L = \lim_{n\to\infty}(x_n)$. Trivially, we have that $\lim_{n\to\infty}(x_{n + 1}) = L = \lim_{n\to\infty}(x_n)$, and thus:

	\begin{align*}
		\lim_{n\to\infty}x_{n + 1} &= \lim_{n\to\infty}\sqrt{2 + x_n} \qquad \langle \; \text{By definition of} \; (x_n) \; \rangle \\
		\lim_{n\to\infty}x_{n + 1} &= \sqrt{2 + \lim_{n\to\infty}x_n} \qquad \langle \; \text{By Algebraic Limit Theorems} \; \rangle \\
		L &= \sqrt{2 + L} \\
		L^2 &= 2 + L \\
		L^2 - L - 2 &= 0 \\
		(L + 1)(L - 2) &= 0 \\
	\end{align*}

	Thus, the limit of $(x_n)$ must be one of either $L = -1$ or $L = 2$. However, since $(x_n)$ is an increasing sequence, and $x_1 = 1$, it naturally must be that $\lim_{n\to\infty}(x_n) \ge 1$. Thus, $\lim_{n\to\infty}(x_n) = 2$.
\end{solution}


















\end{document}
